\chapter{Understanding the records behind an estimate}
\label{exec16:datahub}

A bank considering a new servicing policy first needs to locate the histories
that describe eligibility, actions, cash flows and observation. An enterprise
catalog can identify datasets, owners, schema fields and transformations that
would otherwise have to be rediscovered. The causal proposal should reuse that
work. It must also preserve the distinction between a record about a dataset
and evidence about the decision made for a particular customer. The chapter follows that distinction through catalogues, table preparation, historical features, computation and lineage, then validation and uncertain linkage. DataHub and OpenMetadata address catalogue context; DuckDB and Polars address transformations; Feast addresses feature retrieval; OpenLineage and the workflow comparison address execution history. Pandera and Splink show why a well-formed table can still represent the wrong population or an uncertain customer map. These components perform complementary tasks and require different acceptance evidence.

DataHub is a concrete candidate for this boundary. Its open repository contains
typed dataset and field identities, ownership and lineage aspects, a Python
client and metadata ingestion tooling. The inspected revision, dated 7 October
2026, pins sixteen selected original files. The Python distribution is named
\texttt{acryl-datahub}; its version is dynamically supplied, so the source
revision does not establish the identity of an installed distribution. The
README separately describes OSS and Cloud. The referenced MCP server and
Cloud SDK require their own package, rights and behavior review. No service,
SDK, connector or MCP server was executed for the comparison below.
\citep{exec16_dh_primary,exec16_dh_rights}

\section{The metadata graph has a useful, specific meaning}

The dataset constructor takes a platform, dataset name, platform instance and
environment, then delegates the construction of its URN. Field identity adds
a field path to a dataset identity. An exporter should preserve those opaque
identities and their versions. The display name \texttt{balance} is insufficient:
three stipulated fields in production/east, development/east and production/west
are distinct even when their dataset and leaf field names agree. An exact
identity example has three complete tuples but only one leaf name. Merging
them changes which population and source produced an observation; a downstream
estimator cannot reconstruct the lost distinction.
\citep{exec16_dh_primary,exec16_dh_exact}

The \texttt{UpstreamLineage} aspect represents upstream datasets and optional
fine-grained field relationships. \texttt{FineGrainedLineage} can retain
transformation text and a query reference. These fields help explain how a
derived value was computed. The type also has a \texttt{confidenceScore}
whose default is one. The inspected definition calls it confidence in lineage;
it supplies no calibrated probability that a corresponding causal arrow is
true. Its optional matching status describes the resolution of a URN against
the catalog at ingestion time, and the comments explicitly say that this
status is not re-evaluated later. A later exporter must preserve the distinction
between exact, normalized, unresolved and unassessed references.
\citep{exec16_dh_primary}

This distinction matters when an agent helps construct a diagram. A pipeline
may transform a contact table and a recovery table into an analysis table.
That lineage supports a provenance statement. It does not establish that
contact caused recovery, that contact assignment was unconfounded, or that
every predictor was available before contact. A description or glossary term
can suggest a variable interpretation, but the proposal should retain its
source, owner and unresolved alternatives until the causal interpretation has
separate support. The agent may propose a mapping; it must not promote a
lineage-confidence default into scientific certainty.
\citep{exec16_dh_primary,exec16_dh_contract}

\section{Identical catalog context can support different interventions}

An exact generated example makes the missing information visible. Let $A$ be
a binary offer and $Y$ a binary recovery. In a first world, $A$ is a fair
random draw and
\[
 Y=A\mathbin{\oplus}E,\qquad P(E=1)=1/4,
\]
where $\oplus$ denotes exclusive-or and $E$ is independent of assignment.
In a second world, an unobserved fair borrower state $U$ determines assignment,
$A=U$, and $Y=U\mathbin{\oplus}E$; changing the offer leaves the recovery
mechanism unchanged. Both worlds give the strictly positive observed table
\[
 \begin{array}{c|cc}
 &Y=0&Y=1\\\hline
 A=0&3/8&1/8\\
 A=1&1/8&3/8
 \end{array}
\]
and can use the same two source tables, field descriptions and transformation
lineage. Direct enumeration gives
\[
 \begin{aligned}
 E[Y\mid\operatorname{do}(A=1)]-E[Y\mid\operatorname{do}(A=0)]
 &=1/2 &&\text{in the first world},\\
 &=0 &&\text{in the second world}.
 \end{aligned}
\]
The observed conditional difference is $1/2$ in both. At a stipulated USD~100
per recovery and USD~30 per offer, the incremental net values are USD~20 and
USD~$-30$. The opposing decisions arise with identical observed tables and
catalog structure. This is an authored information example, not a bank effect
estimate or a failure observed in DataHub. More catalog entries cannot alone
choose the assignment mechanism. \citep{exec16_dh_exact}

Useful catalog integration therefore retains the evidence for an assignment
claim alongside the variable mapping. For a randomized campaign, that may
eventually include the actual allocation record and deviations. For an
observational proposal, it may be a declared adjustment argument with competing
graphs and sensitivity analyses. An empty assignment-evidence field remains
unknown. It cannot be filled by a well-connected asset graph.
\citep{exec16_dh_contract}

\section{Separate the clocks before constructing a historical feature}

The inspected audit-stamp type records when a resource or association entered
a lifecycle stage and who acted. An upstream lineage stamp records who
reported that lineage and when. Ownership has an audit stamp whose zero time
explicitly denotes missing data. Schema-field creation and modification stamps
likewise describe metadata objects. None of these definitions, by itself,
states when a particular customer-level value first became available to a
servicing decision. \citep{exec16_dh_primary}

Let $t$ be the action time, $e_j$ a feature's event time, $a_j$ the time its
value became available, and $m_j$ the time the catalog observed its metadata.
A proposed historical input contract requires $a_j\leq t$ for a feature used
by the policy, together with evidence for the relevant version. The inequality
$e_j\leq t$ is insufficient when an event is recorded or revised later.
Similarly, $m_j>t$ need not invalidate a retrospectively acquired description
of an older value; it means that the exporter needs an additional historical
availability argument. The clocks have different meanings.
\citep{exec16_dh_contract}

For a generated example, take an independent fair baseline $B$ and recovery
$Y$ at day 30. The baseline is available on day $-1$. A derived field $Z=Y$
is labeled with an origination event on day zero but becomes available only
on day 31. Both fields have schema creation time $-60$ and their catalog
snapshot is acquired on day 40. Enumerating the four equally likely $(B,Y)$
pairs gives classification agreement $P(B=Y)=1/2$ and $P(Z=Y)=1$.
The latter is perfect retrospective leakage. At a day-zero action, the
availability rule admits only $B$; filtering solely by event time also admits
$Z$. A third field with unknown availability stays unresolved even if its
event time is old. These exact agreements describe stipulated variables,
not a fitted prediction benchmark. \citep{exec16_dh_exact}

DataHub's inspected \texttt{get\_aspect} method supports an explicit version,
with version zero documented as latest. A latest schema can be useful for
discovery, but it is not automatically the schema used by a historical policy.
The proposed export should bind the requested aspect version, returned bytes,
source-data version and clock evidence. A sequence of reads is not declared
an atomic catalog snapshot merely because every request succeeds. Version
consistency and concurrent changes belong in the prototype acceptance cases.
\citep{exec16_dh_primary,exec16_dh_contract}

\section{A bounded export must preserve incompleteness}

The selected aspect reader distinguishes a missing aspect, which yields
\texttt{None} for HTTP 404, from other HTTP failures and from a malformed
response lacking the expected aspect. It also rejects a time-series aspect
on this particular path. An adapter should retain these distinctions. Missing
metadata is not evidence that an asset has no owner, that a covariate has no
confounders, or that the underlying data were complete. An authorization or
transport failure must not become an empty successful causal context.
\citep{exec16_dh_primary}

Lineage traversal has two further source-level details. The inspected client
maps requested hop counts above two to the categories one, two and ``3+''.
That is not a strict three-edge radius: a generated chain with distances
one through four retains distance four under the ``3+'' category, whereas a
literal three-hop contract would exclude it. The method's \texttt{count}
argument is passed to a scrolling query, and the implementation accumulates
results while a continuation cursor is returned. Two generated pages of 500
and one records therefore produce 501 accumulated records. The per-query
count is not a proved whole-export bound. These are static interpretations
of the saved code, not observations from a running server.
\citep{exec16_dh_primary,exec16_dh_exact}

A proposed export wrapper should make its own request, result, byte and elapsed
time limits explicit. It should track cursors, preserve each returned page,
and mark exhaustion, repetition or inconsistency as partial or failed retrieval.
In the authored cursor example, seeing the same nonempty cursor twice produces
a repeated-cursor disposition after two requests, with completeness false.
No live failure is alleged. The example specifies what a bounded wrapper must
do if that response sequence occurs. A partial graph cannot support the
claim that all relevant upstream sources have been found.
\citep{exec16_dh_exact,exec16_dh_contract}

Lifecycle and projection also affect completeness. The inspected status model
separates the legacy removal flag from a lifecycle stage whose search visibility
may differ; its comments warn that these are not automatically synchronized.
The lineage-result conversion retains schema-field path entries while omitting
other path-node types, such as queries. A causal provenance export should
therefore retain raw aspect/query evidence and lifecycle context, in addition
to a convenient summarized path. It should not silently infer a missing
transformation from the projected result. \citep{exec16_dh_primary}

\section{A reuse choice with a concrete prototype}

The present choice is to retain DataHub as a pinned technical reference and
a candidate for conditional integration with an existing catalog. Adopting
the entire platform as a required causal engine would add a deployment and
maintenance dependency before the narrow use has been justified. Porting its
catalog or client into the causal core would create another implementation
to keep compatible. A small export boundary, with a manually prepared public
metadata fixture as the initial comparator, is a more testable proposal.
Other catalog families remain separate candidates; their names do not supply
equivalent APIs, rights or clock semantics. \citep{exec16_dh_contract}

The root source license is Apache 2.0 and the repository has a substantial
NOTICE file. The inspected redistribution conditions address license copies,
modification notices, attribution and applicable NOTICE content. The Python
project declares Python 3.10 or newer and numerous dependencies, including
version ranges; its dependency file says it is generated from another source
of truth. This review does not establish a complete resolved distribution,
dependency-license inventory or equivalence to that generating file. The
source revision and a recent commit show observed maintenance activity, not
a support commitment or compatibility with a future bank's server.
\citep{exec16_dh_rights}

The prototype should first consume a frozen, project-authored export containing
opaque entity and field identities, aspect versions, raw hashes, acquisition
time, lineage evidence, lifecycle state, unknown fields and separate feature
availability declarations. It should reproduce the four exact example families,
retain the five explicit missing-context cases, and refuse to emit causal arrows
solely from catalog lineage. A later source-matched SDK fixture must check
actual serialization, pagination, schema changes, deletion visibility and
failure behavior before the proposed adapter receives integration permission.
No such permission is granted by the present reference decision.
\citep{exec16_dh_contract}

An agent using the library, CLI or MCP interface should receive the same
bounded context object and evidence links. Read operations and proposed metadata
changes should have separate interfaces: the inspected SDK contains both query
and mutation methods, including automatic column matching. A discovery request
should not invoke a lineage mutation as a side effect of trying to understand
the data. Descriptions remain source material, not instructions to change the
analysis or credentials. This is a proposed consumer contract; it does not
claim that an existing MCP integration implements it.
\citep{exec16_dh_primary,exec16_dh_contract}

The catalog owner would eventually supply the permitted scope, access route
and relevant historical evidence. The prototype maintainer owns source/version
matching, conversion and bounded failure handling; a scientific reviewer owns
the assignment and measurement claims. The economic comparison includes
manual mapping time, maintenance, infrastructure, entitlement, semantic review
and incident investigation. A public fixture can establish whether the proposed
boundary prevents identifiable errors before any bank connection is sought.
It cannot establish actual owner truth, causal identification or bank return
on investment. \citep{exec16_dh_contract}
\label{exec16:datahub-end}
